% =============================================================
% BACKGROUND — max 500 words (MethodsX requirement)
% =============================================================

Change detection is one of the most widely used operations in remote
sensing. Comparing two co-registered images of the same place at
different dates reveals where the surface has changed, and the same
operation serves many fields. It maps urban
expansion~\cite{chen2020spatial}, assesses damage after
disasters~\cite{gupta2019xbd}, tracks agricultural and land-cover
transitions~\cite{yang2021asymmetric}, and monitors
deforestation~\cite{prodes,mapbiomas}. The deep-learning methods behind
these applications~\cite{daudt2018fully}, and their shared challenges,
are the subject of dedicated surveys~\cite{khelifi2020deep,shi2020change}.

Supervised change detection is built on pixel-level labels, and
producing them is the main cost of a dataset. An analyst inspects the
imagery, decides the state of each pixel, and paints a dense mask.
Change makes this harder than single-image labeling in two ways. The
unit is a pair of co-registered dates rather than one image, so the
analyst reasons about both at once, and the same scene can be labeled
against several reference dates, which multiplies the work. Expert time
becomes the binding constraint, and reducing it has driven a range of
label-efficient strategies. Weakly supervised methods replace dense
masks with image-level tags or
scribbles~\cite{lin2016scribblesup,hua2021semantic}. Semi-supervised
methods combine a small labeled set with a large unlabeled
one~\cite{zhou2026decaf}. Active learning chooses which samples the
human should label next~\cite{ren2021survey}. Human-in-the-loop
annotation lets the model and the user take turns, the model proposing a
mask and the user correcting it~\cite{amershi2014power,lenczner2022dial}.
Point supervision replaces dense masks with a small set of
clicks~\cite{Bearman2016}, and recent work brings it to change detection
specifically~\cite{fang2024point}.

iSAGE~\cite{isage} (\emph{Iterative Sparse Annotation Guided by Expert})
is a recent point-based, human-in-the-loop framework for remote-sensing
semantic segmentation. An expert clicks only on the model's confident
mistakes and an error-weighted loss amplifies the gradient at each
click, with no pseudo-labels, conditional random fields, or
foundation-model prompts, and the annotation record is itself the
dataset. On single-image benchmarks it recovered most of the quality of
dense labeling from far less than 1\% of the pixels. iSAGE was designed
for the single-image case, and its label file, click loop, and
error-weighted loss all assume one image per annotation. This article
introduces \textbf{iSAGE-CD}, an adaptation of iSAGE to change
detection. iSAGE-CD is a general annotation tool that applies to any
bi-temporal change task on any co-registered imagery. We describe the
tool, its pair-anchored session format, and the annotation workflow, and
we illustrate them with a working annotation session. The aim is to
document a reusable method rather than to run a controlled experiment,
and a full quantitative evaluation is the subject of a separate
manuscript in preparation.
